% DSAIL RA 申请 CV：第一页英文，第二页中文，左上角语言标识
% 编译：xelatex
\documentclass[10pt]{ctexart}

\usepackage[margin=1.5cm]{geometry}
\usepackage{fancyhdr}
\usepackage{enumitem}
\usepackage{needspace}

\pagestyle{fancy}
\fancyhf{}
\renewcommand{\headrulewidth}{0.4pt}
\setlength{\headheight}{14pt}

\setlist[itemize]{leftmargin=1.2em, itemsep=1pt, topsep=1pt, parsep=0pt}

\newcommand{\sectionhead}[1]{%
  \needspace{3\baselineskip}%
  \vspace{6pt}%
  {\normalsize\bfseries #1}%
  \vspace{2pt}%
  \hrule height 0.5pt
  \vspace{3pt}%
}

\newcommand{\name}[1]{{\Large\bfseries #1}}

\begin{document}

% ============ 第一页：English ============
\fancyhead[L]{\small English}
\fancyhead[R]{\small}

\begin{center}
  \name{Lianxiang Li}
\end{center}

\begin{center}
  llxstupg@163.com \textbar\ +86 175 1637 6491 \textbar\ github.com/LI-PG1
\end{center}

\sectionhead{Education}
\begin{itemize}
  \item \textbf{City University of Hong Kong} — MSc in Electronic Commerce (jointly offered by Dept.\ of Computer Science \& Dept.\ of Information Systems), Fall 2026 -- expected Oct 2027 / Feb 2028
  \item \textbf{Anhui University (Project 211)} — BSc in Applied Statistics, 2021 -- 2025 \textbar\ GPA 80/100
\end{itemize}

\sectionhead{Experience}
\textbf{AI Development Engineer, VirtAItech (OrionX), Beijing} \hfill Mar -- Aug 2026
\begin{itemize}
  \item Delivered 30+ open LLMs (text, VLM, MoE, OCR, image/video) to production inference services; owned vLLM configuration, environment adaptation, and acceptance evaluation.
  \item Built reusable RAG and agent frameworks (embedding + rerank over Milvus, tool-calling orchestration); fine-tuned domain models with LoRA/QLoRA and DeepSpeed.
  \item Multimodal customer service for CITIC Securities: speech transcription (SenseVoice), document/chart field extraction (Qwen2.5-VL), finance knowledge base (73\% auto-resolution, $\sim$1.6s first response).
\end{itemize}

\sectionhead{Selected Projects}
\textbf{Text2SQL Agent with Few-Shot Retrieval} — Ctrip Data Platform \hfill Jul -- Aug 2026
\begin{itemize}
  \item Schema linking and few-shot selection over historical query--SQL pairs; 400+ case golden evaluation set (indicator-QA 91\%, SQL execution 88\%).
\end{itemize}

\textbf{Multi-Tool Agent for Equipment Maintenance} — SAIC Motor, POC \hfill Apr -- Jun 2026
\begin{itemize}
  \item Compared ReAct vs.\ Plan-then-Execute; MCP-style tool registration (7 tools); multi-step completion 74\% $\rightarrow$ 86\%.
\end{itemize}

\textbf{Grounded RAG with Evidence Verification} — Zhejiang Telecom \hfill 2026, during internship
\begin{itemize}
  \item On-premise pipeline (4$\times$A100); hybrid retrieval + cross-encoder rerank + LLM-as-Judge verification loop; top-5 hit rate +32\%.
\end{itemize}

\sectionhead{Research Interest}
I am interested in Generative AI, LLMs, and agents. At work I built RAG and agent systems; I would now like to understand these models more deeply, not just deploy them. I am especially interested in how DSAIL applies generative models to scientific problems such as molecular design, and my statistics background makes me curious about the data side of that work.

\sectionhead{Skills}
Python \textbar\ PyTorch, Transformers, DeepSpeed, PEFT (LoRA/QLoRA) \textbar\ vLLM, Docker, FastAPI \textbar\ LangChain, LangGraph, RAG, agent orchestration \textbar\ Milvus, Chroma, Redis \textbar\ probability, regression, multivariate analysis \textbar\ IELTS 6.5

\newpage

% ============ 第二页：中文 ============
\fancyhead[L]{\small 中文}
\fancyhead[R]{\small}

\begin{center}
  \name{李联翔}
\end{center}

\begin{center}
  邮箱：llxstupg@163.com \textbar\ 电话：+86 175 1637 6491 \textbar\ GitHub：github.com/LI-PG1
\end{center}

\sectionhead{教育背景}
\begin{itemize}
  \item \textbf{香港城市大学} — 电子商务理学硕士（计算机科学系与信息系统系合办），2026 秋 — 预计 2027.10 / 2028.02 毕业
  \item \textbf{安徽大学（211）} — 应用统计学学士，2021 — 2025 \textbar\ GPA 80/100
\end{itemize}

\sectionhead{实习经历}
\textbf{大模型应用开发工程师，北京趋动科技（OrionX）} \hfill 2026.03 -- 2026.08
\begin{itemize}
  \item 交付 30+ 个开源大模型（文本/VLM/MoE/OCR/图像与视频）至生产推理服务，负责 vLLM 配置、环境适配与上线评估。
  \item 搭建可复用 RAG 与 Agent 框架（Milvus 上的 Embedding+Rerank、工具调用编排），用 LoRA/QLoRA + DeepSpeed 微调垂域模型。
  \item 参与中信建投多模态客服：语音转写（SenseVoice）、证件/图表字段抽取（Qwen2.5-VL）、金融知识库（自动解决率 73\%，首响约 1.6s）。
\end{itemize}

\sectionhead{代表项目}
\textbf{Text2SQL 智能体 + Few-Shot 检索} — 携程数据中台 \hfill 2026.07 -- 2026.08
\begin{itemize}
  \item 设计 schema linking 与历史 query--SQL 对 few-shot 检索；构建 400+ 条黄金评估集（指标问答 91\%，SQL 执行 88\%）。
\end{itemize}

\textbf{设备运维多工具智能体} — 上汽集团，POC \hfill 2026.04 -- 2026.06
\begin{itemize}
  \item 对比 ReAct 与 Plan-then-Execute 编排；MCP 式工具注册（7 类工具）；多步任务完成率 74\% $\rightarrow$ 86\%。
\end{itemize}

\textbf{带证据校验的落地 RAG} — 浙江电信 \hfill 2026（实习期间）
\begin{itemize}
  \item 私有化部署（4$\times$A100）；混合检索 + cross-encoder 精排 + LLM-as-Judge 校验闭环；top-5 命中率 +32\%。
\end{itemize}

\sectionhead{研究兴趣}
我对生成式 AI、大语言模型和智能体感兴趣。工作里搭建过 RAG 和智能体系统；现在想更深入理解这些模型，而不只是部署它们。我尤其关注 DSAIL 如何把生成模型应用到分子设计这类科学问题上，我的统计背景也让我对这项工作里的数据层面很好奇。

\sectionhead{技能}
Python \textbar\ PyTorch、Transformers、DeepSpeed、PEFT（LoRA/QLoRA）\textbar\ vLLM、Docker、FastAPI \textbar\ LangChain、LangGraph、RAG、Agent 编排 \textbar\ Milvus、Chroma、Redis \textbar\ 概率论、回归、多元分析 \textbar\ 雅思 6.5

\end{document}
